% ======================================================================
%  CAPÍTULO 6 — RESULTADOS E DISCUSSÃO
%
%  Lotes de 27, 29 e 31/10/2026 (análise com ANG a 400 Hz, reanalisados em
%  1º e 2/11). As tabelas saem de ensaios/tabelas_monografia.py --desde
%  2026-10-27 (voo mais recente de cada par); os números do texto, das
%  mesmas pastas de resultado e de monografia/verificacao/.
% ======================================================================
\chapter{Resultados e discussão}
\label{cap:resultados}

Este capítulo apresenta os resultados dos ensaios descritos no
Capítulo~\ref{cap:metodos}: os 49 cenários e as 14 configurações de
controlador da Seção~\ref{sec:controladores}, em 398 pares cenário ×
controlador, voados em 27, 29 e 31/10/2026; vale, para cada par, o voo mais
recente. As Seções~\ref{sec:res-integracao} a~\ref{sec:res-vento} discutem o
PID nativo, com os ganhos da Tabela~\ref{tab:ganhos}, e o \textit{backend} de
exemplo, com os mesmos ganhos, e têm duas funções: verificar que o caminho de
integração pelo \codigo{AC\_CustomControl} reproduz o controlador nativo,
condição para que a comparação com os controladores propostos seja válida, e
caracterizar o PID nativo como linha de base quantitativa. As tabelas dessas
seções já trazem as linhas dos controladores propostos, que a
Seção~\ref{sec:res-proposto} discute.

Cada cenário foi voado ao menos uma vez com cada controlador da sua lista;
quando houve repetição, vale o voo mais recente. Ensaios repetidos depois dos
lotes, com os mesmos ganhos, dão a medida da repetibilidade: em dois ensaios
do LQG repetidos em 1º/11 --- o degrau de torque com o braço encurtado e o
degrau de torque com \codigo{SIM\_GYR1\_RND} = 20 ---, o desvio máximo e o
erro quadrático médio se reproduziram a menos de 5\,\%, com desvios de
\ang{0.433} e \ang{0.413} no primeiro e de \ang{0.522} e \ang{0.525} no
segundo. Com o ruído de giroscópio mais alto, a variação entre voos idênticos
cresce, e a Seção~\ref{sec:res-proposto} a quantifica. A única fonte de
variação entre voos identificada na entrada é a duração efetiva do pulso de
torque, discutida na Seção~\ref{sec:res-torque}.

% ----------------------------------------------------------------------
\section{Verificação do caminho de integração}
\label{sec:res-integracao}

Em todos os ensaios com o \textit{backend} de exemplo, a mensagem
\codigo{CC} do registro mostra que ele assumiu o controle cerca de
\SI{0.4}{\second} depois do início do perfil e permaneceu ativo até o fim,
antes, portanto, de qualquer degrau ou perturbação. Com os mesmos ganhos, os
dois controladores produzem respostas praticamente iguais à perturbação de
torque (Tabela~\ref{tab:res-torque}): no degrau sustentado, os desvios
máximos diferem em no máximo 2\,\% e os erros quadráticos médios coincidem;
no pulso de rolagem e arfagem, os desvios diferem em 1,5\,\%. O modelo da
Seção~\ref{sec:validacao-tempo} reproduz cada controlador separadamente, com
a mesma precisão, a partir do código de cada um.

Duas diferenças aparecem, e ambas têm explicação. A primeira é o
sobressinal do degrau de rolagem e arfagem, maior no \textit{backend}
(Tabela~\ref{tab:res-degrau}): 27\,\% e 24\,\%, contra 22\,\% e 20\,\% no
nativo. A causa é estrutural. O controlador de ângulo nativo é o controlador
de raiz quadrada~\eqref{eq:sqrt-controller}, que, com o ganho de
\SI{26.7}{\per\second} e o limite de aceleração da rolagem, só é linear para
erros de até \ang{1.0}; o do \textit{backend} é puramente proporcional. No
degrau, o erro de ângulo chega a \ang{3.3}, e, nesse erro, o nativo pede uma
velocidade angular 28\,\% menor que a do proporcional, o que reduz o
sobressinal. Em guinada, o limite linear é de \ang{10.9}, o erro do degrau
não passa de \ang{7.9}, e os dois controladores coincidem. A segunda é o
desvio no pulso de guinada: medido, é 6\,\% menor no \textit{backend}, mas o
impulso aplicado no voo dele foi 7\,\% menor (Seção~\ref{sec:res-torque}).
A diferença não vem do controlador: simulados com a mesma entrada, os dois
dão o mesmo desvio, e, normalizados pelo impulso, os desvios são de
\ang{5.55} e \ang{5.64} (Tabela~\ref{tab:res-torque}), 1,6\,\% de
diferença.

O caminho de integração está, portanto, verificado: um controlador executado
dentro do \codigo{AC\_CustomControl} enfrenta a mesma planta, a mesma
perturbação e a mesma análise que o nativo, e as diferenças medidas entre os
dois correspondem às diferenças de código. Para a comparação com os
controladores propostos, fica a lição de que o \textit{backend} não herda o
controlador de raiz quadrada do nativo: o que os controladores propostos
fizerem com erros grandes é escolha de projeto
(Capítulo~\ref{cap:projeto}), e não um detalhe de implementação.

% ----------------------------------------------------------------------
\section{Resposta ao degrau}
\label{sec:res-degrau}

\begin{table}[htbp]
\centering
\caption{Resposta ao degrau de atitude: \ang{10} em rolagem e arfagem e
\ang{30} em guinada.}
\label{tab:res-degrau}
\small
% Gerada por ensaios/tabelas_monografia.py — nao editar a mao.
\begin{tabular}{llrrrrr}
\toprule
Eixo & Controlador & Subida & Sobressinal & Acomodação & \multicolumn{2}{c}{Erro de rastreamento} \\
\cmidrule(lr){6-7}
 & & (\si{\second}) & (\%) & (\si{\second}) & máx. (\si{\degree}) & eficaz (\si{\degree}) \\
\midrule
Rolagem & nativo & 0,045 & 22,1 & 0,18 & 3,25 & 0,36 \\
Rolagem & \codigo{cc\_pid} & 0,040 & 27,1 & 0,23 & 3,31 & 0,39 \\
Rolagem & LQI & 0,045 & 5,6 & 0,10 & 0,59 & 0,06 \\
Rolagem & LQR+ESO & 0,043 & 6,2 & 0,10 & 0,65 & 0,06 \\
Rolagem & LQG & 0,045 & 5,2 & 0,18 & 0,50 & 0,06 \\
Rolagem & LQG (mesmo ruído) & 0,045 & 6,0 & 0,19 & 0,58 & 0,07 \\
Arfagem & nativo & 0,048 & 19,8 & 0,18 & 3,18 & 0,35 \\
Arfagem & \codigo{cc\_pid} & 0,045 & 23,8 & 0,18 & 3,25 & 0,37 \\
Arfagem & LQI & 0,045 & 3,3 & 0,10 & 0,57 & 0,06 \\
Arfagem & LQR+ESO & 0,045 & 3,6 & 0,10 & 0,60 & 0,06 \\
Arfagem & LQG & 0,048 & 5,1 & 0,19 & 0,51 & 0,06 \\
Arfagem & LQG (mesmo ruído) & 0,048 & 5,3 & 0,19 & 0,53 & 0,07 \\
Guinada & nativo & 0,465 & 10,8 & 1,53 & 7,88 & 2,20 \\
Guinada & \codigo{cc\_pid} & 0,465 & 11,0 & 1,54 & 7,88 & 2,19 \\
Guinada & LQI & 0,560 & 0,1 & 0,81 & 0,09 & 0,04 \\
Guinada & LQR+ESO & 0,560 & 0,0 & 0,82 & 0,24 & 0,08 \\
Guinada & LQI (mesmo ruído) & 0,560 & 0,1 & 0,81 & 0,10 & 0,05 \\
Guinada & LQG & 0,568 & 0,3 & 0,82 & 0,11 & 0,04 \\
Guinada & LQG (mesmo ruído) & 0,568 & 0,3 & 0,82 & 0,09 & 0,03 \\
\bottomrule
\end{tabular}


\fonte{Elaborada pelo autor.}
\end{table}

A Tabela~\ref{tab:res-degrau} reúne as métricas dos degraus. Com o nativo e o
\textit{backend} de exemplo, em rolagem e
arfagem, o degrau de \ang{10} sobe em 40 a \SI{48}{\milli\second}, com
sobressinal de 20\,\% a 27\,\% e acomodação em 0,18 a \SI{0.23}{\second}; a
aeronave chega a ficar \ang{3.3} atrás do alvo modelado. Os ganhos de ângulo
do AutoTune são altos --- cerca de seis vezes os do arquivo de parâmetros
padrão --- e o atuador do SITL não tem dinâmica (Seção~\ref{sec:motor-sitl}),
o que permite respostas tão rápidas; numa aeronave real, a dinâmica dos
motores limitaria essa banda. A arfagem repete a rolagem, como prevê a
simetria do quadro X. Em guinada, o degrau de \ang{30} sobe em
\SI{0.46}{\second}, com 11\,\% de sobressinal e acomodação em
\SI{1.5}{\second}, e o erro de rastreamento chega a \ang{8}: o alvo é
modelado com a aceleração máxima de \SI{115}{\degree\per\second\squared}, e
a eficácia de controle em guinada é 30 vezes menor que em rolagem.

% ----------------------------------------------------------------------
\section{Rejeição de perturbações de torque}
\label{sec:res-torque}

\begin{table}[htbp]
\centering
\caption{Rejeição de torque: pulso de \SI{15}{\radian\per\second\squared}
(rolagem e arfagem) e \SI{3}{\radian\per\second\squared} (guinada) por
\SI{0.3}{\second}, e degrau de 10 e \SI{1}{\radian\per\second\squared}; no
pulso de guinada, desvio normalizado ao impulso nominal ($\dagger$).}
\label{tab:res-torque}
\small
% Gerada por ensaios/tabelas_monografia.py — nao editar a mao.
\begin{tabular}{lllrrrr}
\toprule
Eixo & Torque & Controlador & Desvio máx. (\si{\degree}) & Recuperação (\si{\second}) & RMSE (\si{\degree}) & $u_{\max}$ \\
\midrule
Rolagem & pulso & nativo & 0,90 & 1,32 & 0,229 & 0,051 \\
Rolagem & pulso & \codigo{cc\_pid} & 0,91 & 1,21 & 0,224 & 0,051 \\
Rolagem & pulso & LQI & 0,98 & 0,75 & 0,267 & 0,051 \\
Rolagem & pulso & LQR+ESO & 0,99 & 0,73 & 0,275 & 0,051 \\
Rolagem & pulso & LQG & 0,81 & 0,51 & 0,180 & 0,053 \\
Rolagem & pulso & LQG (mesmo ruído) & 0,79 & 0,22 & 0,189 & 0,051 \\
Arfagem & pulso & nativo & 0,93 & 1,25 & 0,233 & 0,050 \\
Arfagem & pulso & \codigo{cc\_pid} & 0,94 & 1,25 & 0,233 & 0,050 \\
Arfagem & pulso & LQI & 1,04 & 0,87 & 0,290 & 0,049 \\
Arfagem & pulso & LQR+ESO & 1,04 & 0,78 & 0,295 & 0,050 \\
Arfagem & pulso & LQG & 0,84 & 0,18 & 0,189 & 0,051 \\
Arfagem & pulso & LQG (mesmo ruído) & 0,86 & 0,17 & 0,198 & 0,051 \\
Guinada & pulso & nativo & 5,55$^\dagger$ & 0,35 & 1,678 & 0,254 \\
Guinada & pulso & \codigo{cc\_pid} & 5,64$^\dagger$ & 0,35 & 1,563 & 0,242 \\
Guinada & pulso & LQI & 3,96$^\dagger$ & 1,68 & 1,402 & 0,225 \\
Guinada & pulso & LQR+ESO & 3,97$^\dagger$ & 1,55 & 1,379 & 0,235 \\
Guinada & pulso & LQI (mesmo ruído) & 7,53$^\dagger$ & 3,25 & 3,477 & 0,132 \\
Guinada & pulso & LQG & 5,35$^\dagger$ & 1,64 & 1,708 & 0,276 \\
Guinada & pulso & LQG (mesmo ruído) & 4,66$^\dagger$ & 1,72 & 1,357 & 0,271 \\
\addlinespace
Rolagem & degrau & nativo & 0,60 & 2,25 & 0,070 & 0,034 \\
Rolagem & degrau & \codigo{cc\_pid} & 0,61 & 2,09 & 0,070 & 0,034 \\
Rolagem & degrau & LQI & 0,65 & 0,44 & 0,049 & 0,035 \\
Rolagem & degrau & LQR+ESO & 0,66 & 0,41 & 0,051 & 0,033 \\
Rolagem & degrau & LQG & 0,51 & 0,17 & 0,031 & 0,034 \\
Rolagem & degrau & LQG (mesmo ruído) & 0,54 & 0,20 & 0,033 & 0,035 \\
Arfagem & degrau & nativo & 0,61 & 2,15 & 0,073 & 0,033 \\
Arfagem & degrau & \codigo{cc\_pid} & 0,62 & 2,22 & 0,073 & 0,033 \\
Arfagem & degrau & LQI & 0,69 & 0,60 & 0,060 & 0,033 \\
Arfagem & degrau & LQR+ESO & 0,69 & 0,46 & 0,056 & 0,033 \\
Arfagem & degrau & LQG & 0,56 & 0,17 & 0,033 & 0,035 \\
Arfagem & degrau & LQG (mesmo ruído) & 0,57 & 0,18 & 0,034 & 0,036 \\
Guinada & degrau & nativo & 2,56 & 22,18 & 0,847 & 0,093 \\
Guinada & degrau & \codigo{cc\_pid} & 2,56 & 22,09 & 0,847 & 0,093 \\
Guinada & degrau & LQI & 1,87 & 1,13 & 0,214 & 0,084 \\
Guinada & degrau & LQR+ESO & 1,77 & 1,02 & 0,193 & 0,087 \\
Guinada & degrau & LQI (mesmo ruído) & 6,16 & 2,33 & 0,987 & 0,082 \\
Guinada & degrau & LQG & 2,28 & 1,34 & 0,261 & 0,097 \\
Guinada & degrau & LQG (mesmo ruído) & 1,90 & 1,61 & 0,230 & 0,093 \\
\bottomrule
\multicolumn{7}{l}{\footnotesize $^\dagger$ Pulso de guinada: desvio máximo normalizado ao impulso nominal do cenário, $\delta_{\max}\,I_{\mathrm{nom}}/I_{\mathrm{reg}}$.} \\
\end{tabular}


\fonte{Elaborada pelo autor.}
\end{table}

A Tabela~\ref{tab:res-torque} mostra a rejeição de torque. Com o nativo e o
\textit{backend} de exemplo, em rolagem e
arfagem, o pulso desvia a atitude em 0,90 a \ang{0.94}, e o erro volta a
10\,\% do pico em 1,2 a \SI{1.3}{\second}; o degrau sustentado de
\SI{10}{\radian\per\second\squared} desvia \ang{0.6} e é compensado pelo
integrador em 2,1 a \SI{2.3}{\second}. O esforço de controle é pequeno: o
pulso pede cerca de 5\,\% da faixa de esforço normalizado, e o degrau,
3,4\,\%. Em guinada, o quadro muda. Um torque sustentado dez vezes menor,
\SI{1}{\radian\per\second\squared}, desvia o rumo em \ang{2.56}, quatro vezes
mais que o de rolagem, e a recuperação leva \SI{22}{\second}, dez vezes mais.
No pulso de guinada, a recuperação de \SI{0.35}{\second} é medida até
10\,\% do pico, mas resta um desvio de cerca de 4\,\% do pico, que o
integrador lento só elimina em dezenas de segundos. A Seção~\ref{sec:validacao-tempo} mostrou que essa recuperação é governada
pelo polo em $-k_I/k_P$, imposto pela razão $k_I = k_P/10$ que o AutoTune
adota em guinada, e não pela planta. É nesse eixo, com a menor eficácia de
controle e a ação integral mais lenta, que a rejeição de perturbação do PID
nativo é mais fraca, e onde um controlador com estimação explícita da
perturbação tem mais a oferecer.

O pulso de torque exige uma ressalva. O laço do injetor roda a cerca de
\SI{44}{\hertz}, e não aos \SI{50}{\hertz} pedidos, e a duração efetiva do
pulso nominal de \SI{0.3}{\second} variou de 300 a \SI{338}{\milli\second}
nos voos de 27/10 do nativo e do \textit{backend} de exemplo (de 300 a
\SI{344}{\milli\second} em todos os voos válidos), e o impulso aplicado, em
até 10\,\% entre dois voos. Em rolagem e
arfagem, a malha fechada é rápida o bastante para que o pico do desvio ocorra
durante o próprio pulso, e ele depende pouco da duração: dois voos com
impulsos 7\,\% diferentes desviaram 0,90 e \ang{0.91}. Em guinada, a resposta
é lenta, e o desvio acompanha o impulso. Por isso, nas tabelas deste
capítulo, o desvio do pulso de guinada é normalizado ao impulso nominal do
cenário, $\delta_{\max}\,I_{\mathrm{nom}}/I_{\mathrm{reg}}$ (marca
$\dagger$); em rolagem e arfagem, em que o pico precede o fim do pulso, vale o
desvio medido. A alternativa seria o injetor temporizar o pulso pelo próprio
simulador, que é capaz de aplicar o torque por uma duração exata
(\codigo{SIM\_TWIST\_TIME}).

\begin{table}[htbp]
\centering
\caption{Rejeição de torque com o PID nativo, no ensaio de cada eixo e no
ensaio simultâneo dos três eixos.}
\label{tab:res-simultaneo}
\small
% Gerada por ensaios/tabelas_monografia.py — nao editar a mao.
\begin{tabular}{llrrrr}
\toprule
 & & \multicolumn{2}{c}{Desvio máx. (\si{\degree})} & \multicolumn{2}{c}{Recuperação (\si{\second})} \\
\cmidrule(lr){3-4}\cmidrule(lr){5-6}
Eixo & Torque & isolado & 3 eixos & isolado & 3 eixos \\
\midrule
Rolagem & pulso & 0,90 & 0,90 & 1,32 & 1,29 \\
Arfagem & pulso & 0,93 & 0,92 & 1,25 & 1,31 \\
Guinada & pulso & 5,55$^\dagger$ & 5,58$^\dagger$ & 0,35 & 0,35 \\
Rolagem & degrau & 0,60 & 0,59 & 2,25 & 2,29 \\
Arfagem & degrau & 0,61 & 0,61 & 2,15 & 2,27 \\
Guinada & degrau & 2,56 & 2,54 & 22,18 & 22,14 \\
\bottomrule
\multicolumn{6}{l}{\footnotesize $^\dagger$ Pulso de guinada: desvio máximo normalizado ao impulso nominal do cenário, $\delta_{\max}\,I_{\mathrm{nom}}/I_{\mathrm{reg}}$.} \\
\end{tabular}


\fonte{Elaborada pelo autor.}
\end{table}

Os ensaios simultâneos nos três eixos (Tabela~\ref{tab:res-simultaneo})
mostram que o acoplamento entre os eixos é desprezível no SITL --- em parte
porque o simulador não modela o termo giroscópico
$\vetor{\omega}\times\matriz{J}\,\vetor{\omega}$ (Seção~\ref{sec:sitl}),
que numa aeronave real acoplaria rolagem e arfagem em velocidades angulares
altas. Com rolagem,
arfagem e guinada perturbadas ao mesmo tempo, e com os motores M3 e M4
trabalhando o dobro do que trabalham num ensaio de um eixo, cada eixo desvia
o mesmo que no ensaio isolado, com diferença de no máximo 1,1\,\%, e
recupera com diferença de no máximo 6\,\%. É o que prevê a matriz de entrada
diagonal da Seção~\ref{sec:modelo-linear}, e justifica projetar o
controlador eixo a eixo.

% ----------------------------------------------------------------------
\section{Sensibilidade à perturbação na entrada da planta}
\label{sec:res-frequencia}

\begin{table}[htbp]
\centering
\caption{Resposta em frequência do torque para o ângulo, nas varreduras de
torque.}
\label{tab:res-frequencia}
\small
% Gerada por ensaios/tabelas_monografia.py — nao editar a mao.
\begin{tabular}{llrrrr}
\toprule
Eixo & Controlador & Ganho máx. & $f$ do máx. (\si{\hertz}) & Queda de \SI{3}{\decibel} (\si{\hertz}) & Coerência \\
\midrule
Rolagem & nativo & 0,063 & 3,6 & 6,8 & 0,99 \\
Rolagem & \codigo{cc\_pid} & 0,063 & 3,6 & 7,2 & 0,99 \\
Rolagem & LQI & 0,084 & 1,6 & 4,8 & 0,99 \\
Rolagem & LQR+ESO & 0,086 & 1,6 & 4,4 & 0,99 \\
Rolagem & LQG & 0,087 & 2,0 & 4,4 & 0,98 \\
Rolagem & LQG (mesmo ruído) & 0,077 & 2,8 & 5,2 & 0,98 \\
Arfagem & nativo & 0,064 & 3,6 & 6,8 & 0,99 \\
Arfagem & \codigo{cc\_pid} & 0,065 & 3,6 & 6,8 & 0,99 \\
Arfagem & LQI & 0,085 & 1,6 & 4,8 & 0,99 \\
Arfagem & LQR+ESO & 0,089 & 1,6 & 4,0 & 0,99 \\
Arfagem & LQG & 0,091 & 2,0 & 4,0 & 0,98 \\
Arfagem & LQG (mesmo ruído) & 0,077 & 2,8 & 5,2 & 0,98 \\
Guinada & nativo & 2,562 & 0,2 & 1,2 & 0,98 \\
Guinada & \codigo{cc\_pid} & 2,564 & 0,2 & 1,2 & 0,98 \\
Guinada & LQI & 2,328 & 0,4 & 1,0 & 0,99 \\
Guinada & LQR+ESO & 2,309 & 0,6 & 1,0 & 0,99 \\
Guinada & LQI (mesmo ruído) & 7,795 & 0,2 & 0,6 & 0,99 \\
Guinada & LQG & 2,540 & 0,6 & 1,2 & 0,98 \\
Guinada & LQG (mesmo ruído) & 3,024 & 0,6 & 1,0 & 0,98 \\
\bottomrule
\end{tabular}


\fonte{Elaborada pelo autor.}
\nota{\raggedright Ganho em \si{\degree} por
\si{\radian\per\second\squared}; queda de \SI{3}{\decibel} em relação ao
ganho máximo, acima da frequência do máximo; coerência mediana na faixa da
varredura.}
\end{table}

As varreduras (Tabela~\ref{tab:res-frequencia}) estimam a resposta do ângulo
ao torque de perturbação, $P\,S$, a sensibilidade à perturbação na entrada da
planta (Seção~\ref{sec:bode}), e mostram em que frequências a malha fechada é
mais sensível a ela. Com o nativo e o \textit{backend} de exemplo, em rolagem
e arfagem, o ganho é máximo em \SI{3.6}{\hertz}, com 0,061 a
\SI{0.065}{\degree} por \si{\radian\per\second\squared}, e cai
\SI{3}{\decibel} entre 6,4 e \SI{7.6}{\hertz}; abaixo do máximo, o integrador
atenua a perturbação lenta, e acima dele, a própria inércia da
aeronave~\cite[seção~11.3]{astrom2008}. Em guinada, o ganho é máximo na menor
frequência da varredura, \SI{0.2}{\hertz}, com \SI{2.56}{\degree} por
\si{\radian\per\second\squared} --- quarenta vezes o máximo de rolagem ---, e
cai \SI{3}{\decibel} em \SI{1.2}{\hertz}. A guinada é sensível justamente às
perturbações lentas, que o integrador fraco demora a compensar, coerente com
o degrau sustentado. A coerência mediana é de 0,97 ou mais em todas as
varreduras, e o modelo linear reproduz essas respostas
(Seção~\ref{sec:validacao-frequencia}).

% ----------------------------------------------------------------------
\section{Vento e torque}
\label{sec:res-vento}

O ensaio de rajada confirma, com medida, a divisão de trabalho entre as duas
famílias de perturbação (Seção~\ref{sec:injetor}). Sob uma rajada lateral de
\SI{10}{\metre\per\second}, a aeronave, com a atitude comandada e sem
controle de posição, chegou a \SI{4.2}{\metre\per\second} e derivou
\SI{18.6}{\metre}, mas o erro de rolagem não passou de \ang{0.017}, contra
\ang{0.013} a \ang{0.015} no trecho calmo antes da rajada, e o de arfagem
ficou em \ang{0.015}. O pulso de torque de rolagem desvia a atitude 50 vezes
mais. A guinada é a única exceção visível: o erro de rumo passou de
\ang{0.03} para \ang{0.30} durante a rajada, ainda vinte vezes menos que no
pulso de torque de guinada. O nativo e o \textit{backend} de exemplo dão os
mesmos números. No modelo de multirrotor do SITL, o vento é, assim,
sobretudo uma perturbação de posição: o arrasto de corpo age no centro de
massa, e o arrasto de momento nos rotores só produz momento na medida em que
os empuxos diferem, o que explica o pequeno efeito em guinada; os ensaios de
rejeição de perturbação da malha de atitude usam, por isso, torque.

% ----------------------------------------------------------------------
\section{Comparação com o controlador proposto}
\label{sec:res-proposto}

\pendente{Comparar o controlador proposto (Capítulo~\ref{cap:projeto}) com o PID nativo nos
mesmos dezesseis cenários, pelas métricas da Seção~\ref{sec:metricas}, com
atenção especial à guinada sob torque sustentado e normalizando os pulsos de
guinada pelo impulso registrado.}
